% results.tex -- Results
\section{Results}

To determine whether a lasting change in one local response could move
a stable market path to a Sell majority, we activated Rule~19286 at
four fractions of interior sites and compared each switched path with
a matched no-switch control, and we found that only the switched paths
crossed (Fig.~\ref{fig:rule-quench}). The selected campaign used a
$64\times64$ lattice, an unlimited dealer, and an endogenous boundary
(Supplementary Table~\ref{tab:campaigns}). Each switched path shared
its 1,000-step history and its first post-activation quote with its
control, so the two paths differed only in how the targeted traders
responded to a neighbor sum of one. All 50 controls remained below a
majority for 3,000 post-activation steps at every fraction, as did the
switched paths at $f=0.10$ and 0.15, whereas activation at $f=0.20$ or
0.25 produced a crossing in every seed, with median onset falling from
1,044.5 to 311.5 steps between the two larger fractions
(Fig.~\ref{fig:rule-quench}B; Supplementary Table~\ref{tab:paired-switch}).
In the representative $f=0.20$ pair, Sell-rich regions expanded through
the switched lattice while the control remained mixed
(Fig.~\ref{fig:rule-quench}C). Two controls asked which part of the
intervention carried this effect. A matched one-time Sell pulse, which
placed the same target sites in Sell once and left their rule
unchanged, produced no majority crossings at $f=0.15$, 0.20, or 0.25
under either boundary, so the lasting response change produced an
outcome that the transient state change did not
(Supplementary Table~\ref{tab:revision-pulse}). An exploratory screen
of all 18 single-entry neighbors of Rule~19280 found majority
transitions for six of them at $f=0.20$ and none for the other twelve
(Supplementary Table~\ref{tab:revision-rules}). Within the tested
interventions, collective selling therefore depended on the persistence
of the response change, on which response changed, and on how many
traders changed it.

A Sell majority could spread through neighbor responses alone or
through the price signal that boundary traders fed back into the
lattice. To separate the two routes, we repeated the switch on a finer
fraction grid across eight dealer capacities under both the endogenous
boundary and a held boundary that kept every active boundary trader in
Hold, and we found that the transition survived removal of the price
signal (Fig.~\ref{fig:reservoir}). Under the held boundary, majority
probability rose from 0.34 at $f=0.175$ to 0.70 at 0.185 and 1.00 at
0.20 (Fig.~\ref{fig:reservoir}A). Changing dealer capacity left every
held-boundary $(f,\text{seed})$ lattice trajectory exactly unchanged,
despite differences in executions, prices, and dealer holdings, because
nothing the dealer did could reach the interior once the boundary
stopped responding to price. Local neighbor responses were therefore
sufficient to sustain collective selling. Restoring the boundary
response to price changed the probability only near the activation
window (Fig.~\ref{fig:reservoir}A,B). With an endogenous boundary and
unlimited dealer, 14 of 50 seeds crossed at $f=0.175$, 43 crossed at
0.185, and all 50 crossed at every sampled fraction from 0.195 upward,
and across all dealer capacities the endogenous runs produced no
transitions at $f\leq0.15$ and transitions in every seed at
$f\geq0.20$. Price feedback thus shifted how often the majority formed
inside the activation window but was not required for it to develop.
The instability lived in the lattice, which left two questions: how
the lattice transition depended on size and time, and what determined
its price consequence.

Because the majority grew out of finite Sell-rich regions, its location
on the activation axis could depend on the size of the lattice and on
how long each run was observed. We repeated the held-boundary switch at
five lattice sizes on refined fraction grids and found that larger
lattices reached a majority at lower activation fractions and over a
narrower crossover (Fig.~\ref{fig:finite-size}). At 3,000 steps under
the held boundary, the midpoint fell from $f_{50}=0.216$ at
$L=32$ to 0.166 at $L=128$, and the width decreased from
0.023 [0.018, 0.027] to 0.008 [0.004, 0.010], where
brackets give 95\% seed-bootstrap intervals
(Fig.~\ref{fig:finite-size}A,C; Supplementary Table~\ref{tab:finite-size-scaling}).
Peak susceptibility, the seed-to-seed variance of the final-window Sell
fraction scaled by the number of interior sites, increased from
21.2 to 241.2 over the
same range (Fig.~\ref{fig:finite-size}B,D). Refined grids resolved each
crossover, and logistic fits under a monotone curve assumption placed
each midpoint within 0.004 of its interpolated value, while the fitted
widths were narrower under the held boundary
(Supplementary Table~\ref{tab:revision-fits}).
Observation time also mattered: extending the $L=64$ window from 3,000
to 9,000 steps lowered the majority midpoint from about 0.178 to 0.167,
and size-scaled windows and alternative Sell-fraction criteria retained
the endpoint size ordering while longer runs continued to shift the
midpoint (Supplementary Table~\ref{tab:revision-horizons}). The sharper
response on larger lattices therefore described a finite-time crossover
whose location depended on the observation horizon, and we did not
treat it as an asymptotic threshold.


Ensemble probabilities did not show how an individual lattice reached a
majority, or whether it did so by gradual spatial growth or by a
near-simultaneous change of many states. We examined the prespecified
time-sliced runs across the activation crossover, which were selected
from ensemble summaries before any image was inspected, and found that
every crossing proceeded by growth and merger of Sell-rich regions
(Fig.~\ref{fig:dose-panels}). At $f=0.15$, the noncrossing run retained
a fine mixed pattern for all 3,000 steps. The two $f=0.175$ runs
followed different paths: one formed a Sell-rich region but stayed
below one-half, while the other developed a large connected domain and
crossed after 2,186 steps. The $f=0.185$ run crossed after 1,779 steps
through several connected regions, and the $f=0.20$ run crossed after
1,030 steps and ended with one dominant Sell domain
(Fig.~\ref{fig:dose-panels}A). Price declined as Sell occupancy rose in
each crossing run, with maximum drawdowns of 17.6\%, 23.7\%, and
25.6\% at $f=0.175$, 0.185, and 0.20, against 5.3\% and 9.9\% for the
two noncrossing runs (Fig.~\ref{fig:dose-panels}B,C). The number,
shape, and timing of the domains varied across seeds, and the
noncrossing $f=0.175$ run also formed a visible domain, so the
appearance of a Sell-rich region did not by itself guarantee a majority
within the horizon. Each displayed transition therefore proceeded by growth and merger of
Sell-rich domains rather than by a synchronous relabeling of the
lattice, and the ensemble probability counted how often a majority was
reached within the horizon.

Having located the instability in the lattice, we asked what determined
its price consequence. We compared majority probability, maximum
drawdown, and dealer rejection across the eight dealer capacities under
the endogenous boundary and found that capacity ordered the price
decline while changing majority probability non-monotonically
(Fig.~\ref{fig:reservoir}). Near the activation window, at $f=0.175$,
majority probability was 0.66 with $h=0.1$ and 0.28 with unlimited
capacity, a 38-percentage-point advantage for the smaller dealer
(Fig.~\ref{fig:reservoir}A,B). Freezing each path's boundary decision
width at its pre-switch value changed this contrast from 0.38 to
$-0.02$, which implicated the adaptation of that width to the moving
quote in the advantage while leaving other non-monotone differences in
place (Supplementary Table~\ref{tab:revision-band}). At $f=0.20$, every
seed transitioned under every capacity, but greater capacity allowed
more eligible orders to execute and so increased the drawdown: median
maximum drawdown rose from zero at $h=0$ to 0.40\% at $h=0.1$, 3.92\%
at $h=1$, and 25.3\% with unlimited capacity, while median dealer
rejection fell from 0.449 to zero (Fig.~\ref{fig:reservoir}C,D). Zero
capacity admitted only balanced trader flow and left the quote fixed,
and $h=10$ rejected no orders and matched unlimited capacity for these
endowments and horizon. The ordering of drawdowns by capacity survived
a decade-wide sweep of the price-impact coefficient, although the
magnitudes changed by a factor of seven to ten across the sweep
(Supplementary Table~\ref{tab:revision-impact}). The same separation
appeared within one seed. At $f=0.20$, seed 25 reached a majority at
$h=0.1$, 1, and $\infty$ within 88 steps of one another, and its
final-window mean Sell fraction lay between 0.81 and 0.87 in all three
endogenous paths and in the held-boundary path, yet its maximum
drawdown rose from 0.40\% at $h=0.1$ to 3.92\% at $h=1$ and 25.8\%
with unlimited capacity (Fig.~\ref{fig:dealer-panels}). A dealer could
therefore limit the price decline after most interior traders had
entered Sell, while price feedback changed how often the majority
formed near the window.

Dealer capacity limited which eligible orders executed, but trader
holdings determined which orders were eligible in the first place. To
separate the two resources, we multiplied trader cash and shares
independently by 0.1, 1, or 10 at activation under three dealer
capacities and both boundaries, and we found that holdings changed
majority probability only when prices fed back to the lattice, whereas
their effect on drawdown depended on the dealer
(Fig.~\ref{fig:balance-sheet}). At $f=0.18$, all nine cash--share
interventions produced identical held-boundary lattice paths and a
majority probability of 0.64 across dealer capacities, while endogenous
probabilities ranged from 0.50 to 0.94, and the other sampled fractions
showed the same boundary dependence (Fig.~\ref{fig:balance-sheet}A;
Supplementary Table~\ref{tab:balance-fractions}). Holdings altered order
eligibility, but their effect on local behavior required the resulting
price signal. With baseline cash and an unlimited dealer, increasing
the share multiplier from 0.1 to 1 to 10 raised median drawdown from
2.87\% to 20.1\% to 83.5\%, because every eligible sell executed and
the dealer absorbed the whole unmatched residual
(Fig.~\ref{fig:balance-sheet}B). Finite
capacity capped that increase: baseline and tenfold shares both
produced 3.92\% at $h=1$ and 0.399\% at $h=0.1$, because the dealer
rejected the unmatched sells that exceeded its capacity
(Fig.~\ref{fig:balance-sheet}C). Majority probability did not follow
drawdown (Fig.~\ref{fig:balance-sheet}D). At $h=1$ and baseline shares,
tenfold cash reduced median drawdown from 3.92\% to 0.089\% while
raising majority probability from 0.78 to 0.94, because the extra cash
made more buy orders eligible, and those buys matched sells before they
reached the dealer while the resulting quote also changed the boundary
signal. Trader and dealer resources therefore determined how much
collective selling translated into a price decline, and a Sell majority
and a large drawdown were separate outcomes of the same intervention.
